\Ch{Pipeline Stages}{Stages}

\p Shader programs are arranged by runtime drivers into a series of steps called
a \textit{pipeline}. Depending on the type of workload being performed the
pipeline may contain shaders implementing different \textit{stages}.

\p The \textit{basic compute} stage is called the \textit{compute} stage.

\p The \textit{basic rasterization} stages are the \textit{vertex} and
\textit{fragment} stages.

\begin{note}
A fragment shader may also be called a \textit{pixel} shader.
\end{note}

\p The \textit{mesh primitive} stages are \textit{amplification} and
\textit{mesh}.

\p The \textit{ray tracing} stages are \textit{ray generation},
\textit{intersection}, \textit{any hit}, \textit{closest hit}, \textit{miss},
and \textit{callable}.

\p During shader compilation the target shader stage must be know to the
compiler. This may be provided by the compiler from the user, or with a source
attribute (\ref{Decl.Attribute.Shader}). Declarations defined in the standard
library clauses (\ref{Library}) may restrict availability based on target shader
stage. A program is ill-formed if it relies on a declaration that is not
available in the target shader stage. An implementation may omit diagnostics for
provably unreachable.

\p All conforming implementations must support the basic compute stage; all
other stages are optional except that any conforming implementation that
supports the mesh primitive stages or ray tracing stages must also support the
basic rasterization stages.

\p Each shader stage defines a set of special named input and output variables
called \textit{system values}. A stage defines which system values are legal as
inputs or outputs for each stage and what the values mean.

\Sec{Compute Stage}{Stages.Compute}

\p A compute shader represents a generic computation and may be used for general
purpose programmability.

\Sub{Input System Values}{Stages.Compute.Input}

\Sec{Basic Rasterization Stages}{Stages.Rasterization}

\p Each \gls{lane} in a vertex shader produces a mutated vertex value from a set
of stage-provided and user-defined inputs. Each \gls{lane} in a fragment shader
produces a color value for a rendered fragment which contributes to the final
color value for one or more output pixels.

\Sub{Input System Values}{Stages.Compute.Input}
\Sub{Output System Values}{Stages.Compute.Output}

\Sec{Mesh Primitive Stages}{Stages.Mesh}

\p A mesh shading pipeline may optionally contain an amplification shader which
is responsible for dispatching \gls{threadgroup}s and creating payloads for the
mesh shader to operate on. A mesh shader operates on a small batch of vertices
and primitives grouped together called a \textit{meshlet}, and performs a
similar role to a vertex shader.

\Sub{Input System Values}{Stages.Compute.Input}
\Sub{Output System Values}{Stages.Compute.Output}

\Sec{Ray Tracing Stages}{Stages.RayTracing}

\p A ray tracing shader pipeline is comprised of a ray generation shader which
begins the dispatch and launches rays into the scene. The pipeline may register
intersection, closest hit, and any hit shaders to execute on ray collision
conditions, or miss shaders to execute when no collision occurs. Callable
shaders may be directly invoked from other ray tracing shader stages.

\Sub{Input System Values}{Stages.Compute.Input}
\Sub{Output System Values}{Stages.Compute.Output}
